\section{Introduction}
A hard per-image symbol budget must accommodate source information, signaling and protection. We focus on preserving source structure and perceptual similarity, primarily evaluated by DISTS and DreamSim, while retaining PSNR for pixelwise fidelity. Transmitted information constrains generative reconstruction to the source; a learned prior predicts untransmitted content. The description's organization and use of payload capacity therefore affect reconstruction quality.

Next-scale tokenization provides complete coarse maps over the full latent grid, followed by residual refinements~\cite{var}. Within a scale, token probabilities share the completed coarse context, enabling scale-wise evaluation and completion. ARPC applies this representation to progressive compression~\cite{arpc}; Ada-TokenCom combines next-token prefix coding, tail generation and wireless adaptation~\cite{adatokencom}. These systems motivate studying next-scale descriptions under hard digital-frame budgets.

Progressive transmission has long matched source descriptions to channel resources~\cite{chande}. A full next-scale map may exceed remaining capacity although much of it would fit. Image-dependent entropy lengths also prevent a fixed token count from specifying the bit requirement. Using residual capacity requires independently decodable within-scale endpoints and actual-length fitting. Received boundary tokens must then update both the residual representation and finer-scale conditioning.

NeST-Com combines token-wise endpoints, independently terminated arithmetic coding and prefix-constrained reconstruction. Shared-model comparisons and matched controls establish improvements in pixelwise fidelity and perceptual similarity. At 1024 symbols and 10 dB, source-bit utilization rises from 77.70\% to 99.79\%. Equal-information-budget tests identify a favorable tight-budget regime, while matched-hardware measurements show shorter processing than the fixed-frame Ada implementation. These results distinguish the practical value of the next-scale system from the contribution of its frame-fitting design.
